\chapter{Kurzfassung}

\begin{german}
	Die Generierung realistischer Landschaften für Anwendungen wie Videospiele ist oft
	durch die mangelnde Kontrolle über geografische Details begrenzt. Diese Bachelorarbeit
	untersucht, wie Höhenkarten, die von dem Diffusionsmodell MESA \cite{borne-pons2025mesa} generiert werden,
	präziser gesteuert werden können, um geografisch akkurate Landschaften zu erzeugen.
	Eine Evaluierung von ControlNet \cite{zhang2023controlnet} für die Integration in MESA wurde durchgeführt. Das Training und die Bewertung erfolgen auf dem Major\_TOM Datensatz \cite{schulz2024majortom} unter Verwendung von algorithmisch extrahierten Kontrollkarten. Ziel ist es, durch zusätzliche Kontrollsignale eine verbesserte Steuerbarkeit und Genauigkeit der
	generierten Landschaften zu erreichen, die über die textbasierte Steuerung hinausgeht.
\end{german}
